\documentclass[main.tex]{subfiles}
\begin{document}
% ============================================================
%  1. Motivating Examples   (~12 min)
% ============================================================
\section{Motivating Examples}

% ------------------------------------------------------------
\subsection{An elliptic model problem}
\begin{frame}{Example 1: A Simple Inverse Problem}
  \framesubtitle{See, e.g., Lions (1971), Tr\"oltzsch (2010)}
  \visible<2->{
  \begin{qbox}
   Given the measured temperature $u_d$ in a room $D$. Reproduce the heat source $z$.
  \end{qbox}
  }  

  \visible<3->{
  \begin{hlbox}
   For simplicity, assume the stationary temperature $u$ in $D$ solves
     \[\aligned
    -\Delta u &= z \;\;\text{in } D,\qquad
            u = 0 \;\;\text{on } \partial D,
    \endaligned
  \]
  It is reasonable to assume that $z$ is not arbitrarily large or small. Take $z \in Z_{\rm ad}$ with
  \[
    Z_{\rm ad} := \{ z \in \Ltwo : z_a \le z \le z_b \text{ a.e.\ in } D \}.
  \]
  \end{hlbox}
  }  

  \visible<4->{
  \begin{hlbox}
  Find $z \in Z_{\rm ad}$ such that $u$ is close to $u_d$. An objective function could be:
        \[
    J(u,z) := \frac12 \norm{u - u_d}_{\Ltwo}^2 + \frac{\alpha}{2}\norm{z}_{\Ltwo}^2
  \]
  \end{hlbox}
  }
\end{frame}

% ------------------------------------------------------------
\subsection{A time-dependent problem}
\begin{frame}{Example 2: Optimal Control of the Heat Equation}
  \visible<2->{
  \begin{qbox}
  Suppose I would like to heat the room over a time interval $[0,T]$. How should I heat it to reach a final temperature profile $u_T$?
  \end{qbox}
  }  
  \visible<3->{
  \begin{hlbox}
  The PDE changes to reflect the dynamics ($Q_{T} := D \times (0,T)$),
  \[
    \partial_t u - \Delta u = z \;\;\text{in } Q, \qquad u = 0 \;\;\text{on } \partial D \times(0,T), \qquad u(0) = u_0.
  \]
  We now have $Z_{\rm ad} := \{ z \in L^2(Q_T) : z_a \le z \le z_b \text{ a.e.\ in } Q_{T} \}$.
    \end{hlbox}
  }
  \visible<4->{
  \begin{hlbox}
  The objective is now to find $z \in Z_{\rm ad}$ such that $u(T)$ is close to $u_{T}$
        \[
    J(u,z) :=  \frac{1}{2}\norm{u(T) - u_T}_{\Ltwo}^2 + \frac{\alpha}{2}\norm{z}_{L^2(Q_{T})}^2
  \]
  \end{hlbox}
}
  % \begin{itemize}
    % \item State space $\WT = \{ y \in L^2(0,T;\Hzero) : \partial_t y \in L^2(0,T;H^{-1}(\Omega)) \}$.
    % \item Preview: the adjoint equation runs \alert{backward in time}.
  % \end{itemize}
  % \todo{Insert from your .tex file}
\end{frame}

\begin{frame}{Example 3: Catalyst Placement in a Chemical Reactor}
  \framesubtitle{See, e.g., Tr\"oltzsch (2010), Ch.~4; Casas--Tr\"oltzsch (quasilinear case)}
  \visible<2->{
  \begin{qbox}
   A catalyst density $z$ controls how fast a diffusing reactant $u$ is consumed in a reactor $D$. Where to place $z$ so that $u \approx u_d$.
  \end{qbox}
  }

  \visible<3->{
  \begin{hlbox}
   The stationary concentration $u$ solves the reaction--diffusion equation
   \[
     -\nabla\cdot\big(\kappa(u)\nabla u\big) + z\,R(u) = f \;\;\text{in } \Omega,\qquad
     u = 0 \;\;\text{on } \partial\Omega,
   \]
   feed $f$, diffusivity $\kappa(u) \ge \kappa_0 > 0$, nondecreasing rate law, e.g., $R(u) = \frac{u}{1+u}$ or $u^3$.
  \end{hlbox}
  }

  \visible<4->{
  \begin{hlbox}
  Catalyst bounds and similar objective:
  \[
    Z_{\rm ad} := \{ 0 \le z_a \le z \le z_b \text{ a.e.} \}, \qquad
    J(u,z) := \tfrac12 \norm{u - u_d}_{\Ltwo}^2 + \tfrac{\alpha}{2}\norm{z}_{\Ltwo}^2 .
  \]
  \end{hlbox}
  }

  % \visible<5->{
  % \begin{qbox}
  % \textbf{What changes?} $z \mapsto u$ is \emph{nonlinear} $\Rightarrow$ the problem is \emph{nonconvex}.
  % \end{qbox}
  % }
\end{frame}

% ------------------------------------------------------------
\subsection{A problem under uncertainty}
\begin{frame}{Example 4: Optimization under Uncertainty}
  \visible<2->{
  \begin{qbox}
  Suppose the diffusivity $\kappa(u)$ has been estimated from data and is only understandable through realizations of a random field $\kappa : D \times \Omega \to \mathbb R$. What changes?
  \end{qbox}
  }

\visible<3->{
\begin{hlbox}
    We now view the PDE \textbf{almost surely} wrt to a probability measure $\mathbb P$:
      \[
    -\nabla\cdot\big(\kappa(\cdot,\omega)\nabla u(\omega)\big) + z\,R(u) = f \;\;\text{in } D, \qquad u(\omega) = 0 \;\;\text{on } \partial D, \quad \mathbb P-\text{a.s.}
  \]
\end{hlbox}
}

\visible<4->{
\begin{hlbox}
    While $z$ is still deterministic, $u$ is now a random field in $L^p(\Omega,\mathcal{F},\mathbb P; H^1_0(D))$.
\end{hlbox}
}

\visible<5->{
\begin{hlbox}
The objective $J(u(\omega),z)$ is now implicitly \textbf{random}. We scalarize with a \textbf{risk measure} $\mathcal{R}$ (mean, average-value-at-risk, ...):
  \[
J(u,z) := \mathcal{R}\Big[\tfrac12 \norm{u(\cdot) - u_d}_{\Ltwo}^2\Big] + \frac{\alpha}{2}\norm{z}_{\Ltwo}^2
  \]
\end{hlbox}
}
\end{frame}

% ------------------------------------------------------------
\subsection{Summary}
\begin{frame}{Summary: PDE-Constrained Optimization}
  \visible<2->{
  \begin{hlbox}
  All four examples are instances of
  \[
    \min_{(u,z)}\; J(u,z) \quad\text{subject to}\quad e(u,z) = 0, \qquad z \in Z_{\rm ad}.
  \]
  \end{hlbox}
  }

  \visible<3->{
  \begin{hlbox}
  \begin{itemize}
    \item $z$: \textbf{control}, design, source \hfill (decision variable)
    \item $u$: \textbf{state} \hfill (determined by $z$ through the PDE $e(u,z) = 0$)
    \item $J$: \textbf{objective} \hfill (tracking, terminal observation, risk, \dots)
    \item $Z_{\rm ad}$: \textbf{admissible controls} \hfill (bounds, physical limits)
  \end{itemize}
  \end{hlbox}
  }

  \visible<4->{
  \begin{qbox}
  If the PDE has a unique solution $u = S(z)$ for each $z$, we can eliminate the state:
  \(
    \min_{z \in Z_{\rm ad}} j(z) := J(S(z),z).
  \) This is called the \textbf{reduced formulation}.
  \end{qbox}
  }
\end{frame}

\begin{frame}{Common Structure}
  \begin{center}
  \small
  \renewcommand{\arraystretch}{1.35}
  \begin{tabular}{@{}lllll@{}}
    \toprule
    & 1: Poisson & 2: Heat & 3: Catalyst & 4: Uncertainty \\
    \midrule
    State space & $\Hzero$ & $\WT$ & $\Hzero$ & $L^p_{\mathbb{P}}(\Omega;\Hzero)$ \\
    Control space & $\Ltwo$ & $L^2(Q_T)$ & $\Ltwo$ & $\Ltwo$ \\
    Map $z \mapsto u$ & linear & affine & nonlinear & nonlinear, per $\omega$ \\
    Objective & tracking & terminal & tracking & risk of tracking \\
    Adjoint & elliptic & backward parabolic & linearized elliptic & one per $\omega$ \\
    Reduced problem & convex & convex & nonconvex & nonconvex \\
    \bottomrule
  \end{tabular}
  \end{center}
  \vspace{1ex}
  \visible<2->{
  \begin{qbox}
   Does an optimal control exist? How do we characterize it? How do we compute it?
  \end{qbox}
  }
\end{frame}

\end{document}
